\chapter{RESULTADOS E DISCUSSÕES}
\label{capitulo-resultados-discussoes}

% TODO: preencher na fase de relatório final de extensão (após a execução do projeto)
% Os resultados podem ser apresentados em forma de tabelas ou gráficos, sendo numerados sequencialmente e discutidos antes de serem colocados. Quando possível os resultados experimentais obtidos devem ser comparados com dados de literatura e suas diferenças (quando houver) discutidas.
% Compõe-se de: escrita do assunto a ser tratado no relatório, relato do fato, informação que se queira, descrição de norma e do procedimento. Apresentar somente assuntos que tenham relação direta e específica com o trabalho.

